\chapter{形成性研究}\label{chap:formative}

为厘清历史档案重建场景下专家对计算辅助系统的真实需求，本研究在系统原型设计前开展了一轮形成性研究（formative study）。其目的不在于验证已有界面，而是通过半结构化访谈识别历史学者面对大规模婚姻边补全任务时的认知负担与方法论顾虑，进而抽象出可指导后续视图与协商协议设计的任务谱与设计目标。

\section{研究方法}\label{sec:formative-method}

形成性研究采用半结构化访谈作为主要工具，辅以任务驱动的现场推演。访谈对象通过滚雪球抽样招募，先联系清史人口学研究网络中的初始受访者，再请其推荐方法学背景互补的同行。每轮访谈时长约 60--75 分钟，在线进行，征得书面同意后录音并人工转录；转录稿先由受访者审阅删除可识别身份的片段，再做主题编码。

访谈问题模板由三部分组成。其一为开放式提问，邀请受访者描述其在使用中国多代人口面板（China Multi-Generational Panel Dataset，CMGPD-LN）\citep{campbell2011data} 等历史档案时遇到的婚姻字段缺失问题及应对策略；其二为任务驱动环节，研究者向受访者展示一份 1882 年队列的婚姻候选清单（约 200 对），请其口述如何识别可信对、处理边界对，并为方法学审阅留出可追溯线索；其三为系统假想推演，研究者依次描述若干交互机制（批量接受、协商日志展开、跨队列规则迁移等），请受访者评判其方法学可接受性与潜在风险。整个流程参考混合主动设计\citep{green2009mixedinitiative} 与基于知识生成的可视化分析框架\citep{sacha2014knowledge}。

\section{访谈对象与背景}\label{sec:formative-experts}

本研究共访谈 5 位领域专家，编号 E1--E5，姓名与单位均不披露，仅以角色描述其背景。E1 为清史人口学者，长期关注 1749--1909 年辽宁旗人婚姻策略，具有 5 年以上 CMGPD 研究经验，熟悉八旗户口黄册字段结构\citep{shiue2022marriage}；E2 为计算族谱学者，研究方向偏数字人文中的家谱图重建，熟悉父系世系图的拓扑特征与谱学符号习惯\citep{zhao2001chinese}；E3 为数字人文方法论研究者，关注档案不确定性的形式化表达，对模型可解释性与交互式机器学习\citep{fails2003interactiveml, amershi2014modeltracker} 有理论贡献；E4 为长期使用 CMGPD 的历史学博士生，以辽东南某区域为案例撰写学位论文，熟悉数据清洗与时间窗划分的实务难点\citep{shan2025reexamining}；E5 为民俗学与亲属制度研究者，关注通婚圈、母家与旗内婚的社会规范，对档案中身份用语的细微差异有职业敏感度\citep{duff2003list}。受访者在史学训练与计算方法亲和度上呈互补分布，足以揭示不同立场的方法论顾虑。

\section{关键发现与任务抽象}\label{sec:formative-tasks}

通过对转录稿的开放编码与主题归并，研究者抽象出三项贯穿所有受访者的核心任务，下文以 T1--T3 标识，并以概括叙述方式给出对应的访谈印证。

\textbf{T1：高置信批量接受。} E1 与 E4 均强调，面对数百对候选清单，他们不希望逐对审阅；其偏好是先在宏观层次看清整队列分布，再针对置信度极高的子集做批量决定。E1 描述的理想流程是先观察队列在嵌入空间的聚集形态，识别出与已确认婚姻紧邻的高密度区域，再一键接受证据最充分的若干对。E4 补充，批量接受亦是对模型成熟度的隐性测试——若批次中出现明显反例，便会回退到逐对模式。该偏好印证了状态保持型可视分析对宏观概览的需求\citep{endert2017state}。

\textbf{T2：边界案例评议。} E2、E3、E5 一致认为，真正消耗专家精力的是中等置信、需要结合史料判断的边界对。E5 特别提及，旗内婚与跨旗婚在档案描述上往往只差一个身份用语，模型若仅给出标量分数难以支持决策，专家需要看到结构证据、生活时间线与协商分歧。E2 则关注当模型与史学常识冲突时，系统是否提供足够的反例视角。共同诉求指向一种在模型置信与档案语境之间往返切换的工作流。

\textbf{T3：累积决策审计。} E3 作为方法论研究者，强调任何由计算辅助系统接受的婚姻边都必须可被独立复核：来源是图嵌入、多智能体协商，还是两者协同？E4 出于学位论文规范同样关心可追溯性，期望当审稿人质疑某条决策时能即时回放当时的模型分数、协商日志与专家干预记录。E5 补充，审计需求不应损害交互流畅度，否则专家会绕过审计直接编辑。综合诉求是端到端可追溯且不可篡改的累积决策审计能力。

\section{设计目标推演}\label{sec:formative-dg}

将 T1--T3 与受访者反复提及的方法论顾虑对照后，本研究归纳出六项设计目标 DG1--DG6，作为后续系统设计的硬约束。

\textbf{DG1：拆解图学习为结构证据。} 由 T2 中 E2、E5 对结构可见性的诉求催生，要求系统将异构图变换器（Heterogeneous Graph Transformer，HGT）所捕捉的隐含拓扑模式抽取为可命名的子图模式（micro-motif），并以自然语言陈述供智能体在协商中引用，从而避免标量分数遮蔽证据来源。

\textbf{DG2：全过程审计跟踪。} 直接源于 T3 与 E3、E4 的方法论需求，要求系统对每一条被接受的婚姻边附加来源标记（provenance chip），并完整保留多轮协商日志与人工干预记录，使任何已接受的决策都可在事后被独立复核。

\textbf{DG3：推理前过滤减少计算成本。} 由 T1 中 E1、E4 的批量决策偏好催生，要求在启动昂贵的多智能体推理前，提供基于嵌入空间聚类与模型置信的空间过滤入口，使专家能先收窄候选集再调度计算资源。

\textbf{DG4：推理中干预注入领域知识。} 由 T2 中 E5 对身份用语差异的关切与 E2 对反例视角的需求共同推动，要求在协商进行中提供命令通道，专家可对特定智能体的论证逻辑加权、削弱或终止；该机制与事后合理化（post-hoc rationalization）风险的告诫\citep{turpin2023faithfulness} 契合，也呼应生成式智能体在结构化角色下的协商范式\citep{park2023generative}。

\textbf{DG5：决策可逆并保持状态一致。} 由 T1 与 T3 共同催生，所有受访者在系统假想推演阶段都明确要求一键回退；E4 特别强调回退后必须使被解除的候选回到候选池且全局状态一致，避免出现界面假性成功而后台仍持有旧状态的隐性裂缝。

\textbf{DG6：通过空间先验跨队列复用规则。} 源于 E1 与 E2 提出的"今年校准的规则明年仍可用"的诉求；要求把专家在某一队列上整定的特征权重与子图模式偏好锚定到嵌入空间的局部区域，并在分析新队列时自动激活，从而将单次校准的方法学投入分摊到后续队列。

DG1--DG6 构成本系统设计契约，下游章节的取舍均可回溯至其中之一。

\section{形成性研究小结}\label{sec:formative-summary}

综合访谈可知，历史档案重建场景下的混合主动设计远不止于"让专家接管模型输出"这一表层目标，而需在三个层次回应方法学需求：其一，把模型置信与史学语境联立呈现，避免任何一方独断；其二，为决策设计可追溯且可逆的状态层；其三，把单次队列上的方法学投入沉淀为可复用的空间先验。形成性研究的贡献在于把抽象的"可解释性"诉求拆解为 T1--T3 三项可设计任务，再落地为 DG1--DG6 六项可验证约束。
